% Part: methods
% Chapter: proofs
% Section: introduction

\documentclass[../../../include/open-logic-section]{subfiles}

\begin{document}

\olfileid[id]{mth}{prf}{int}

\olsection{Pendahuluan}

Berdasarkan pengalaman Anda dalam pengantar logika, Anda mungkin
sudah terbiasa dengan !!a{derivation} system---barangkali sistem
!!{derivation} deduksi alami atau bergaya Fitch, atau mungkin sistem
pohon bukti. Anda mungkin ingat pernah membuat bukti dalam sistem-sistem
tersebut, baik untuk membuktikan !!a{formula} maupun untuk menunjukkan
bahwa suatu argumen tertentu sahih. Untuk melakukannya, Anda
menerapkan kaidah-kaidah sistem sampai memperoleh hasil akhir yang
diinginkan. Ketika bernalar \emph{tentang} logika, kita juga membuktikan
berbagai hal, tetapi dalam kebanyakan kasus kita tidak menggunakan
!!a{derivation} system. Bahkan, sebagian besar bukti yang kita bahas
ditulis dalam bahasa Inggris (barangkali dengan sedikit bahasa
simbolik yang disisipkan), bukan seluruhnya dalam bahasa logika orde
pertama. Ketika menyusun bukti semacam itu, pada mulanya Anda mungkin
bingung---bagaimana saya membuktikan sesuatu tanpa !!a{derivation}
system? Bagaimana saya memulainya? Bagaimana saya tahu bahwa bukti saya
benar?

Sebelum mencoba membuat bukti, penting untuk mengetahui apa itu bukti
dan bagaimana menyusunnya. Sesuai dengan namanya, sebuah \emph{bukti}
dimaksudkan untuk menunjukkan bahwa sesuatu itu benar. Anda dapat
membayangkannya sebagai dialog---seseorang bertanya apakah sesuatu itu
benar, misalnya apakah setiap bilangan prima selain dua merupakan
bilangan ganjil. Menjawab ``ya'' saja tidak cukup; orang itu mungkin
ingin tahu \emph{mengapa}. Dalam hal ini, Anda akan memberinya bukti.

Dalam percakapan sehari-hari, mungkin cukup untuk sekadar memberi
isyarat menuju suatu jawaban atau memberikan jawaban yang tidak
lengkap. Namun, dalam logika dan matematika, kita menginginkan bukti
yang ketat---kita ingin menunjukkan bahwa sesuatu itu benar tanpa
menyisakan \emph{keraguan apa pun}. Ini berarti bahwa setiap langkah
dalam bukti harus dibenarkan dan pembenarannya harus kuat (yaitu,
asumsi yang Anda gunakan memang benar-benar diasumsikan dalam
pernyataan teorema yang sedang dibuktikan, definisi yang Anda gunakan
harus diterapkan dengan benar, pembenaran yang dirujuk harus berupa
inferensi yang benar, dan seterusnya).

Biasanya kita membuktikan suatu pernyataan. Pernyataan yang kita
buktikan disebut dengan berbagai nama: proposisi, teorema, lema, atau
korolari. Proposisi adalah pernyataan dasar yang layak dibuktikan:
cukup penting untuk dicatat, tetapi mungkin tidak terlalu mendalam
atau tidak sering diterapkan. Teorema adalah proposisi yang penting
dan berarti. Buktinya sering diuraikan menjadi beberapa langkah, dan
kadang-kadang dinamai menurut orang yang pertama kali membuktikannya
(misalnya Teorema Cantor dan teorema L\"owenheim--Skolem) atau menurut
fakta yang dibahasnya (misalnya teorema kelengkapan). Lema adalah
proposisi atau teorema yang digunakan dalam pembuktian suatu hasil
yang lebih penting. Agak membingungkan, beberapa lema merupakan hasil
penting dengan sendirinya dan juga dinamai menurut orang yang
memperkenalkannya (misalnya Lema Zorn). Korolari adalah hasil yang
dengan mudah mengikuti hasil lain.

Pernyataan yang hendak dibuktikan sering memuat asumsi-asumsi yang
memperjelas jenis objek yang sedang kita buktikan sesuatu tentangnya.
Pernyataan itu mungkin diawali dengan ``Misalkan $!A$ adalah
!!a{formula} berbentuk $!B \lif !C$'' atau ``Andaikan $\Gamma \Proves
!A$'', atau kalimat serupa. Asumsi-asumsi ini merupakan
\emph{hipotesis} proposisi, teorema, atau lema tersebut, dan Anda boleh
menganggapnya benar dalam bukti. Hipotesis membatasi apa yang sedang
kita buktikan dan juga memperkenalkan beberapa nama untuk objek yang
sedang dibicarakan. Misalnya, jika proposisi Anda diawali dengan
``Misalkan $!A$ adalah !!a{formula} berbentuk $!B \lif !C$'', Anda hanya
membuktikan sesuatu mengenai semua rumus dari jenis tertentu (yaitu
kondisional), dan dipahami bahwa $!B \lif !C$ adalah suatu kondisional
sembarang yang akan dibahas oleh bukti Anda.

\end{document}
